\section{Selecting the signed part of the force}\label{sel:section}

Fix a supported receiver and an interval $J=[t_1,t_2]$ of length $I>0$
on which $w/8\le q_X\le8w$. Throughout this section we assume
\eqref{eq:bootstrap} and
\begin{equation}\label{sel:short}
 I\le (w/P)^2.
\end{equation}
Write $L=\log(2+P)$ and $S(w)=P^2L/w$. The estimates proved above give
\begin{equation}\label{sel:baseline}
 \mathcal B(I,w)=P\sqrt I+(M+\sqrt A)I S(w)
\end{equation}
as a bound, up to a fixed factor, for the direct transport and magnetic-remainder
terms and the remaining-sector good-field term. They also give
\begin{equation}\label{sel:absolute}
 \int_J|K(t)|\,dt\le C\sqrt w\,\mathcal B(I,w).
\end{equation}
The extra factor $\sqrt w$ in this absolute estimate will be canceled by
the small coefficient of the receiver terms after signed integration by
parts.

Only the narrow-angle contribution controlled by $S_g$ remains outside
the baseline bound. We estimate bins directly when a near or far bound
for $S_g$ is already below a summable threshold. On the other bins we
apply the signed identity to the original force kernels, including the
transport term. The main task is to show that these selected bins have
a sufficiently small coefficient on the receiver acceleration.

We use the small-sector representatives $p,\theta,\phi,h$, with
$\theta<\kappa\phi$, $\theta\ge1/p$, $p\le P$, and $\phi,h\le1$.
All constants below may depend on the fixed decomposition and on
\begin{equation}\label{sel:epsilon}
 \epsilon=10^{-5},\qquad \beta=1.16,
\end{equation}
but are independent of $M,A,P,I,w$. A sufficiently large lower threshold
for $P$ may depend on $M,A$.

\subsection{The selection and its summable weight}

Put $m=\min(p,w)$ and recall the small-sector stability length
\begin{equation}\label{sel:lambda}
 \lambda^{-1/2}
 =C_1\frac Pm\left(\frac M\phi+\sqrt{\frac{AL}{\phi}}\right).
\end{equation}
Put $I_m^0=(m/P)^2$. The three occupation bounds available for selection
come from energy, stable source-time cells, and the direction estimate,
respectively. Define $\mathsf M_*$ to be the minimum of the applicable
entries in the following table. The displayed entries contain no implicit
constants.
\begin{equation}\label{sel:pure}
\begin{array}{c|l}
 \text{entry}&\text{condition for inclusion}\\ \hline
 \max(I,h)/p&\text{always}\\[2pt]
 (h/p)\max(1,I/\lambda,h/\lambda)&\phi>C_0/m\\[2pt]
 (h/p)\phi^{-\beta}\max(1,I/I_m^0,h/I_m^0)
 &p,w\ge P^{.7},\quad
   j^{-.49}\le\phi\le j^{-.16}\ (j=p,w),
\end{array}
\end{equation}
The occupation estimates give $\mathsf M\le C\mathsf M_*$ in every small
bin. The distinction between this inequality and the exact entries of
\eqref{sel:pure} is useful: the selection below is an exact inequality
between representative quantities.

We assign a summable weight to each bin:
\begin{equation}\label{sel:weight}
 W=(P\theta)^{-\epsilon}\phi^\epsilon h^\epsilon.
\end{equation}
The same weight has two roles. The good-field contribution of an
unselected bin will be at most $CP\sqrt I\,W$ by direct estimation. For a selected bin, we will prove
that the coefficient of $|K|$ in the receiver terms is at most
$Cw^{-1/2}W$. Summing the latter coefficients will cancel the factor
$\sqrt w$ in \eqref{sel:absolute}. The remaining task will then be to
control the central residual and the cutoff and time-weight derivatives
in the signed formula.

A small bin is \emph{selected} if
\begin{equation}\label{sel:selection}
 \min\left\{I\sqrt h\,p^2\phi^2,
       \frac{\sqrt I\sqrt{p\mathsf M_*}}{\theta h}\right\}
       >P\sqrt I\,W.
\end{equation}
This is a fixed set of indices during $J$. In particular, taking a
partial derivative of a cutoff does not differentiate the selection.

\begin{lemma}\label{sel:weight-sum}
The sum of $W$ over all small-sector indices is bounded by a constant
depending only on $\epsilon$ and the fixed decomposition. The unselected
small-sector good-field terms have total integral at most $CP\sqrt I$.
\end{lemma}

\begin{proof}
The $\theta$ sum satisfies
\[
 \sum_{\theta\ge1/p}(P\theta)^{-\epsilon}
 \le C_\epsilon(p/P)^\epsilon.
\]
The remaining sums in $p\le P$, $\phi\le1$, and $h\le1$ are geometric.
The small sector has no additional $\sigma$ partition.
For an unselected bin, the near and far good-field estimates, together
with $\mathsf M\le C\mathsf M_*$, bound its integral by
$CP\sqrt I\,W$. Summing proves the second assertion.
\end{proof}

\subsection{The receiver coefficient on selected bins}

We first verify the coefficient bound promised by the selection. The
source-time identity produces both the receiver-acceleration term
$\mathcal A$ and a derivative of the receiver-angle cutoff. By
\eqref{can:bulk-bounds} and \eqref{can:receiver-cutoff-bound}, both have
majorant
$Cf\theta|K|/(wr\phi^2)$. At a fixed receiver time, the density bound
$\int_{\rm bin}f\,dv\le Cp^3\theta^2$ and the angular cap give a
coefficient at most $Cw^{-1}h^2(p\theta)^3$ on $|K|$. Alternatively,
the particle cone-flux bound gives $C/(wp\theta h\phi^2)$. Thus the
coefficient, as recorded in \eqref{can:receiver-coefficient}, is at most a fixed multiple of
\begin{equation}\label{sel:U}
 U=w^{-1}\min\left\{h^2(p\theta)^3,
                         (p\theta h\phi^2)^{-1}\right\}.
\end{equation}

\begin{lemma}[Receiver coefficient on selected bins]\label{sel:exponents}
For fixed $M,A$, there is a threshold $P_*(M,A)$ such that every selected
index with $P\ge P_*(M,A)$ satisfies the exact representative inequality
\begin{equation}\label{sel:safe}
 U\le w^{-1/2}W.
\end{equation}
\end{lemma}

\begin{proof}
We argue by contradiction: a selected bin violating \eqref{sel:safe}
will first be forced into the actual regime, then into the
improved-occupation window, where its exponent inequalities are
incompatible.

\noindent\textbf{Step 1: an unsafe selected bin has actual velocity separation.}
Write the representatives as exact powers of $P$:
\begin{equation}\label{sel:powers}
 \begin{gathered}
 p=P^{1-\alpha},\qquad w=P^z,\qquad c=1-z,\\
 p\theta=P^b,\qquad \phi=P^{-y},\qquad h=P^{-H},\qquad
 \Delta=\epsilon(\alpha+b+y+H).
 \end{gathered}
\end{equation}
All terminating decimals below denote exact rational numbers.
Here $\alpha,z,c,b,y,H$ are nonnegative and $z+c=1$. The weight is
$W=P^{-\Delta}$. Near selection and \eqref{sel:short} give
\begin{equation}\label{sel:near-exponent}
 2\alpha+H/2+2y<z+\Delta.
\end{equation}
Far selection with the first pure bound in \eqref{sel:pure} gives
\begin{equation}\label{sel:far-energy}
 \begin{cases}
 b<H/2-\alpha+\Delta,&H\le2c,\\
 b<H-c-\alpha+\Delta,&H>2c.
 \end{cases}
\end{equation}
Indeed, when $H\le2c$ one has $h\ge I$, whereas otherwise
$\max(I,h)\le P^{-2c}$. In either case $b\le H-\alpha+\Delta$.
Putting $d_*=\alpha+b+y+H$, we obtain
\[
 d_*\le2H+y+\Delta<4z+5\Delta.
\]
It follows that
\begin{equation}\label{sel:delta-small}
 z>0,\qquad
 \Delta<\frac4{99995}z<.001z.
\end{equation}
Only the pure energy bound has been used so far.

Suppose \eqref{sel:safe} fails. Both terms of the minimum in
\eqref{sel:U} must exceed $w^{-1/2}W$, so
\begin{equation}\label{sel:unsafe}
 b>\frac{2H}3+\frac z6-\frac\Delta3,
 \qquad b<H+2y-\frac z2+\Delta.
\end{equation}
If $H\le2c$, the first inequalities of \eqref{sel:far-energy} and
\eqref{sel:unsafe} force $H+z+6\alpha<8\Delta<.008z$, which is
impossible. Hence $H>2c$, and the other far bound gives
\begin{equation}\label{sel:H-lower}
 H>\frac z2+3(\alpha+c)-4\Delta
 >.496z+3(\alpha+c).
\end{equation}
Combining this with \eqref{sel:near-exponent},
\begin{equation}\label{sel:actual-margin}
 3.5\alpha+1.5c+2y<.753z,\qquad
 z>\frac{500}{751}>.66,\qquad
 y<.377z-1.75\alpha-.75c.
\end{equation}
In particular,
\[
 z-y>.623z+1.75\alpha+.75c\ge.623,
\]
and
\[
 1-\alpha-y>.623z+.75\alpha+1.75c\ge.623.
\]
Thus $m\phi>P^{.3}$, and the bin is actual for sufficiently large $P$.
This conclusion precedes any absorption of constants into powers of $P$.

\smallskip\noindent\textbf{Step 2: stable-cell occupation forces that bin into the improved window.}
Put $\rho=\max(\alpha,c)$, so $P/m=P^\rho$. Equation
\eqref{sel:lambda} and $\phi\le1$ give
\[
 \lambda^{-1/2}
 \le C_1(M+\sqrt{AL})P^{\rho+y}
 \le P^{\rho+y+.001z}
\]
once $P\ge P_*(M,A)$. The last threshold is uniform over the indices
because $z>.66$. Since $I,h\le P^{-2c}$ and $\rho\ge c$, the actual
pure bound gives
\[
 \sqrt{p\mathsf M_*}
 \le P^{-H/2+\rho-c+y+.001z}.
\]
Far selection therefore implies
\begin{equation}\label{sel:far-actual}
 b<H/2-\min(\alpha,c)+y+.002z.
\end{equation}
Comparing \eqref{sel:far-actual} with the first inequality in
\eqref{sel:unsafe}, then using \eqref{sel:near-exponent}, yields
\begin{equation}\label{sel:middle}
 2y>H/3+.328z+2\min(\alpha,c),\qquad
 H<.8076z<.81z.
\end{equation}
Together with \eqref{sel:H-lower}, this gives
\begin{equation}\label{sel:high-momenta}
 \alpha+c<.105z,\qquad z>\frac{200}{221}>.90,
 \qquad1-\alpha>.895.
\end{equation}
The two inequalities in \eqref{sel:unsafe} also imply
\[
 y>\frac z3-\frac H6-\frac{2\Delta}3>.19z.
\]

We now check both ends of the improved-occupation window for both
momenta. The last inequality gives
$y>.171>.16\max(z,1-\alpha)$. For the opposite end put
$Y=.377z-1.75\alpha-.75c$; by \eqref{sel:actual-margin}, $y<Y$, and
\begin{align*}
 .49z-Y&=.113z+1.75\alpha+.75c>0,\\
 .49(1-\alpha)-Y&=.113z+1.26\alpha+1.24c>0.
\end{align*}
Therefore
\[
 .16\max(z,1-\alpha)<y<.49\min(z,1-\alpha).
\]
Together with \eqref{sel:high-momenta}, these are exactly the conditions
for including the improved entry of \eqref{sel:pure}.

\smallskip\noindent\textbf{Step 3: improved occupation contradicts an unsafe coefficient.}
Since $I_m^0=P^{-2\rho}$, that entry gives
\[
 \sqrt{p\mathsf M_*}
 \le P^{-H/2+\rho-c+\beta y/2}.
\]
Far selection and \eqref{sel:unsafe} consequently give
\[
 b<H/2-\min(\alpha,c)+\beta y/2+\Delta,
 \qquad
 \beta y>H/3+.33z+2\min(\alpha,c).
\]
On the other hand, multiplying \eqref{sel:near-exponent} by
$\beta/2=.58$ and using \eqref{sel:delta-small} yields
\[
 \beta y<.58z-.29H+.58\Delta<.582z-.29H.
\]
Thus $H<(378/935)z<.41z$, contradicting \eqref{sel:H-lower}.
This proves \eqref{sel:safe}.
\end{proof}

\subsection{The spatial transition between selected and unselected bins}

The selected receiver terms now have the coefficient needed to use
\eqref{sel:absolute}. We next estimate the source, spatial, and
time-weight derivatives, together with the central residual
$\mathcal R$, in the signed formula. The spatial transition requires
particular care: its angular sum must retain the joint energy budget
without introducing another logarithmic factor.

Equation~\eqref{can:spatial-cutoff-bound} gives the spatial-cutoff majorant
\begin{equation}\label{sel:spatial-kernel}
 C\frac{f}{r^2}\frac\phi\theta.
\end{equation}
By Lemma~\ref{can:cutoff-cancellation}, the sum of the partial source or spatial derivatives of selected cutoff
products is the negative of the corresponding sum over unselected
products. It is supported where selected and unselected products overlap.
Indeed, the full smooth partition sums to one, and a nonnegative smooth
cutoff has zero derivative where it vanishes. Thus we estimate these
errors on overlapping unselected supports. Neighboring representatives
have ratios bounded by four.

For fixed $p,\phi$, let
\begin{equation}\label{sel:R}
 R=1+\sqrt{I/\lambda}\sqrt{\mathcal E_{p\phi}},
 \qquad \sum_{p,\phi}\mathcal E_{p\phi}\le C.
\end{equation}
Outside the deficit exception $\phi\le C_0/m$ and
$h\le\min(I,\lambda)$, the near and far
bounds for \eqref{sel:spatial-kernel} are
\begin{equation}\label{sel:spatial-near-far}
 C\min\left\{Ihp^3\theta\phi^3,
                 \frac{R^2}{ph\theta\phi}\right\}.
\end{equation}
For $h>I$ the pure energy bound $\mathsf M\le Ch/p$ gives the second
bound because $R\ge1$. Hence $R$ is independent of both $\theta$ and $h$
in every use of \eqref{sel:spatial-near-far}.

\Needspace{6\baselineskip}
\begin{lemma}\label{sel:spatial-sum}
The total spatial-cutoff error is at most $C\mathcal B(I,w)$.
\end{lemma}

\begin{proof}
First consider an unselected structural neighbor, with
$\theta\ge\kappa\phi$, overlapping a selected support. Its ratio
$\theta/\phi$ lies within fixed factors of $\kappa$, so for each $\phi$
there are only a bounded number of such $\theta$. Balancing
\eqref{sel:spatial-near-far} in $h$ gives $C\sqrt I\,Rp\phi$, which fits
the shared-budget calculation following \eqref{dir:small-budget}:
the sum of $C\sqrt I\,Rp\phi$ over $p,\phi$ is at most
$C\mathcal B(I,w)$.

For another overlapping unselected bin, fix $p,\phi,\theta$ and define
$h_0>0$ by
\begin{equation}\label{sel:h0}
 I\sqrt{h_0}\,p^2\phi^2
 =P\sqrt I(P\theta)^{-\epsilon}\phi^\epsilon h_0^\epsilon.
\end{equation}
The ratio of the near quantity to its threshold is
$(h/h_0)^{1/2-\epsilon}$. Overlap with a selected neighbor therefore
implies $h\ge c h_0$, for a fixed $c>0$. If the unselected bin fails the
near inequality, it also has $h\le h_0$. Only a bounded number of radial
dyads occur in this case.

Put
\begin{equation}\label{sel:n-f}
 n_0=Ih_0p^3\theta\phi^3,
 \qquad f_*=\frac{R^2}{ph_0\theta\phi},
 \qquad D_0=(n_0f_*)^{1/2}=\sqrt I\,Rp\phi.
\end{equation}
The near-failing transition bins cost at most $C\min(n_0,f_*)$.
If instead the bin fails the far inequality, then
\[
 \mathsf M\le C\mathsf M_*
 \le CP^2W^2\theta^2h^2/p.
\]
Substitution into the far spatial bound gives, by \eqref{sel:h0}, the
additional estimate
\[
 \frac{CP^2W^2\theta}{p\phi}
 =Cn_0(h/h_0)^{2\epsilon}.
\]
Together with \eqref{sel:spatial-near-far}, this bounds the sum over
$h\ge ch_0$ by
\begin{equation}\label{sel:radial-crossing}
 C\begin{cases}
 f_*,&f_*\le n_0,\\
 n_0^{1/(1+2\epsilon)}f_*^{2\epsilon/(1+2\epsilon)},&f_*>n_0.
 \end{cases}
\end{equation}
To see this, split the dyadic sum of
$\min(n_0t^{2\epsilon},f_*/t)$ at
$t=(f_*/n_0)^{1/(1+2\epsilon)}$, where $t=h/h_0$. Both resulting sums
are geometric. The same bound includes the near-failing transition bins.
No restriction on the size of $h_0$ is needed; enlarging the actual
radial range to $t\ge c$ only increases this upper bound.

Equation \eqref{sel:h0} gives
\[
 h_0\propto\theta^{-\epsilon/(1/2-\epsilon)},
 \qquad n_0\propto\theta^\gamma,
 \qquad \gamma=\frac{1-4\epsilon}{1-2\epsilon}>0.
\]
Since $D_0$ is independent of $\theta$, writing $x=n_0/D_0$ turns
\eqref{sel:radial-crossing} into
\[
 CD_0\begin{cases}
 x^{(1-2\epsilon)/(1+2\epsilon)},&x\le1,\\
 x^{-1},&x\ge1.
 \end{cases}
\]
Both $\theta$ tails are geometric. Their sum is at most $CD_0$, with no
additional logarithm. Summing $C\sqrt I\,Rp\phi$ over $p,\phi$
as above gives $C\mathcal B(I,w)$.

Finally, in the omitted deficit exception one has
$\lambda\le CP^{-2}$ and $\phi\le C/m$. The near spatial bound, summed
over $h$ and $\theta$, is at most $CIp^3\phi^4/P^2$. The remaining
momentum weights are bounded by $1/(P^2p)$ when $p\le w$ and by
$P/w^4$ when $p>w$. These fit $CI S(w)$, including the fixed logarithmic
counts from the decomposition. This treats the exception without using
the unavailable far occupation bound there.
\end{proof}

Source-cutoff derivatives have the direct good-field and magnetic-remainder
majorants in \eqref{can:source-cutoff-bound}. On overlapping unselected small bins, Lemma
\ref{sel:weight-sum} applies to the good-field part; the magnetic-remainder part
and structural neighbors are bounded by \eqref{sel:baseline}. The
central term in \eqref{can:bulk-bounds} has majorant $Cf/r^2$. Its
near/far balance is $C\sqrt I\,Rp\theta$, whose $\theta$ sum is
$C\sqrt I\,Rp\phi$. Its deficit exception is bounded by the same near
estimate used at the end of Lemma \ref{sel:spatial-sum}.

\subsection{Time weights}

Choose $0<\eta<1/2$ and an endpoint-zero Lipschitz function $\psi$ on
$J$, equal to one away from monotone ramps of length $\eta I$, with
$0\le\psi\le1$ and $|\psi_t|\le C/(\eta I)$. Equation~\eqref{can:time-weight-bound}
gives the time-weight majorant $C_\eta f\theta/(Ir\phi)$.

\begin{lemma}\label{sel:time-weight}
The total time-weight derivative is at most $C_\eta\mathcal B(I,w)$.
\end{lemma}

\begin{proof}
If $h\le I$, the ratio of this majorant to $f/r^2$ is bounded by
$C_\eta h\theta/(I\phi)\le C_\eta$. The preceding $f/r^2$ estimate
therefore applies.

If $h>I$, the near integral, divided by $P\sqrt I$, is at most
\begin{align}\label{sel:long-weight}
 \frac{C_\eta h^2p^3\theta^3\phi}{P\sqrt I}
 &\le C_\eta h(p/P)^5\phi^3(P\theta\sqrt h)^3W^{-1}\notag\\
 &\le C_\eta h(p/P)^5\phi^3W^{-4}.
\end{align}
The first inequality is the near-selection inequality. For the second,
use far selection and the pure bound $\mathsf M_*\le h/p$ to obtain
$P\theta\sqrt h<W^{-1}$. Equivalently,
\[
 (P\theta)^{1-\epsilon}
 <\phi^{-\epsilon}h^{-1/2-\epsilon}.
\]
Summing the increasing factor $(P\theta)^{4\epsilon}$ in the last
line of \eqref{sel:long-weight} leaves
\[
 C_\eta(p/P)^5
 h^{1-6\epsilon/(1-\epsilon)}
 \phi^{3-4\epsilon/(1-\epsilon)}.
\]
Both exponents are positive. The sums in $h,\phi,p$ are geometric, so
these terms cost at most $C_\eta P\sqrt I$.
\end{proof}

\subsection{The weighted signed estimate}

\begin{proposition}\label{sel:bulk-impulse}
For the weight $\psi$ above and $P\ge P_*(M,A)$,
\begin{equation}\label{sel:bulk}
 \left|\int_J\psi(t)K(t)\,dt\right|
 \le C_\eta\big[P\sqrt I+(M+\sqrt A)I S(w)\big].
\end{equation}
The constant $C_\eta$ is independent of $M,A$.
\end{proposition}

\begin{proof}
Let $\chi=\sum_{j\in\mathcal S}\beta_j$ be the sum of the selected
cutoffs, with the index set $\mathcal S$ fixed throughout $J$. Split the
interior force density as
\[
 T_{\rm ker}+S_{\rm ker}
 =\chi(T_{\rm ker}+S_{\rm ker})
      +(1-\chi)(T_{\rm ker}+S_{\rm ker}).
\]
Although the selection criterion uses the good-field majorant $S_g$,
the selected term contains the transport kernel as required by
Lemma~\ref{can:identity}. Apply Proposition~\ref{can:weighted-formula}
to that term with $\Pi=\Id$, and use the direct bounds on its
complement. Lemmas \ref{sel:weight-sum},
\ref{sel:spatial-sum}, and \ref{sel:time-weight} bound the unselected
good-field terms and all spatial and time-weight derivatives. The
source derivatives, the central signed term, the unselected transport
and magnetic-remainder terms, and structural neighbors were bounded above by
$C_\eta\mathcal B(I,w)$. By \eqref{can:initial-boundary-bound}, the initial-cone boundary has a bounded-data
cost $CI$, which fits $CP\sqrt I$ by \eqref{sel:short}; the endpoint
weights remove the other time boundaries.

For the receiver terms, Lemmas \ref{sel:weight-sum} and
\ref{sel:exponents} give, at every receiver time,
\[
 \sum_{\mathrm{selected}}CU\le Cw^{-1/2}.
\]
Multiplying this coefficient by \eqref{sel:absolute} bounds their total
integral by $C\mathcal B(I,w)$. All sums used fixed geometric exponents;
their constants are independent of $M,A$. The only enlargement of the
large-$P$ threshold occurred in Lemma \ref{sel:exponents}. This proves
\eqref{sel:bulk}.
\end{proof}
